The reduced revision rate for interactively extracted CalibrationTargets is observational rather than controlled. The batch pipeline used GPT-5.1 and the interactive pipeline used Claude Opus 4.6; these models differ in context handling and characteristic error profiles, so any per-mode difference may partly reflect model capability rather than collaboration mode. The interactive work also followed curation of the batch outputs, meaning the modeler entered interactive extraction with working knowledge of where the LLM tended to fail on this corpus. The data are consistent with interactive curation reducing post-hoc revision, but a controlled comparison would require the same underlying model in both modes and randomized assignment of targets.
